% ===== 卒業論文用LaTeXメインファイル(by HONDA) =====
% 分割して作成しているTeXファイルを1つの文書として生成するためのファイル

% 数式を中央寄せしたい場合は fleqn を削除する
\documentclass[autodetect-engine,dvi=dvipdfmx,ja=standard,a4paper,fleqn,12pt]{bxjsreport}
\usepackage{subfiles} % ファイル分割用のパッケージ
\usepackage[bachelor]{tuats} % ページレイアウトの設定． 基本的に中身は触らないほうが安全
\usepackage{mysettings} % 追加で色々設定をまとめておくためのファイル． 必要に応じて加筆すると良い

% タイトルページ関連
\title{藤田桂英研究室\LaTeX~テンプレート 2023} % 日本語タイトル
\etitle{A \LaTeX~Template Files for Fujita Katsuhide Labratory} % 英語タイトル
\author{苗字 名前} % 名前漢字表記
\eauthor{your name} % 名前英語表記
\enteryear{2022} % 入学年度
\graduateyear{2023} % 卒業年度
\studentnumber{12345678} % 学籍番号
\date{YYYY年MM月DD日} % 提出日

\begin{document}

% ==== タイトルページと目次 ====
\maketitle

% ==== メインページ ====
\setcounter{page}{1}
\pagenumbering{arabic}

% hoge.texを読み込む際は\subfile{hoge}とせよ．拡張子は不要．
\subfile{manual}
\subfile{hoge}

% ==== 謝辞(acknowledgement.tex) ====
\subfile{acknowledgement}

% ==== 参考文献(reference.bib) ====
\bibliography{reference}

% ==== 付録 ====
% chapterの前後を付録用に変更．
% 普通は必要無いが，\titleformatでchapterの体裁を書き換えたら「第A章」のようになってしまったので，力技で修正．これで問題なく表示される．
\appendix
\renewcommand{\prechaptername}{付録}
\renewcommand{\postchaptername}{}

% 以下にTeXファイルを読み込んだ場合は「付録」となる．
\subfile{hoge}

\end{document}